% ID: SEL-D0082-OB-E02
% Tanda S082.
\textbf{a)} Cada tramo es continuo en su intervalo. La continuidad en $6$ exige $6a-36=12$, de donde $\boxed{a=8}$.

\textbf{b)} Para $a=8$, en $0<t<6$ se cumple $B'(t)=8-2t$, que cambia de positivo a negativo en $4$. En $6<t<10$, $B'(t)=2>0$. Por tanto, crece en $[0,4]$, decrece en $[4,6]$ y vuelve a crecer en $[6,10]$.

La gráfica es el arco de parábola $8t-t^2$ desde $(0,0)$ hasta $(6,12)$, con vértice $(4,16)$, unido al segmento de recta $B=2t$ hasta $(10,20)$. Los puntos de unión pertenecen a la función.
\begin{center}
\setlength{\unitlength}{1mm}
\begin{picture}(119,75)(-10,-10)
\put(-4,0){\vector(1,0){109}}\put(0,-5){\vector(0,1){65}}
\put(107,0){$t$}\put(1,60){$B$}
\qbezier(0,0)(30,60)(60,30)
\qbezier(60,30)(80,40)(100,50)
\put(0,0){\circle*{1.4}}
\put(40,40){\circle*{1.4}}
\put(60,30){\circle*{1.4}}
\put(100,50){\circle*{1.4}}
\put(1,-6){\scriptsize $(0,0)$}
\put(29,44){\scriptsize $(4,16)$}
\put(59,24){\scriptsize $(6,12)$}
\put(86,54){\scriptsize $(10,20)$}
\end{picture}
\end{center}

\textbf{c)} Restringidos a los primeros seis años, el cambio de signo de la derivada y la comparación $B(0)=0$, $B(4)=16$, $B(6)=12$ muestran que el máximo es $\boxed{16}$ millones de euros a los $\boxed{4}$ años.
